\Chapter{22}

\Verse{1}
{
    \zA{话说贾琏听凤姐儿说有话商量，因止步问：“什么话？” 凤姐道：“二十一是 薛妹妹的生日，你到底怎么样？” 贾琏道：“我知道怎么样？ 你连多少大生日都料
        理过了，这会子倒没有主意了！” 凤姐道：“大生日是有一定的则例。 如今他这生 日，大又不是，小又不是，所以和你商量。” 贾琏听了，低头想了半日，道：“你
        竟糊涂了。 现有比例，那林妹妹就是例。 往年怎么给林妹妹做的，如今也照样给薛 妹妹做就是了。”}
}
{
    \eA{Chia Lien, for we must now prosecute our story, upon hearing lady Feng
        observe that she had something to consult about with him, felt constrained
        to halt and to inquire what it was about. "On the 21st," lady Feng
        explained, "is cousin Hsüeh's birthday, and what do you, after all, purpose
        doing?" "Do I know what to do?" exclaimed Chia Lien;}
}
{
}

\Verse{2}
{
    \zA{凤姐听了冷笑道：“我难道这个也不知道!我也这么想来着。 但 昨日听见老太太说，问起大家的年纪生日来，听见薛大妹妹今年十五岁，虽不算是
        整生日，也算得将笄的年分儿了。 老太太说要替他做生日，自然和往年给林妹妹做 的不同了。” 贾琏道：“这么着，就比林妹妹的多增些。” 凤姐道：“我也这么想
        着，所以讨你的口气儿。 我私自添了，你又怪我不回明白了你了。” 贾琏笑道：“罢! 罢!这空头情我不领。 你不盘察我就够了，我还怪你？”}
}
{
    \eA{"you have made, time and again, arrangements for ever so many birthdays of
        grown-up people, and do you, really, find yourself on this occasion without
        any resources?" "Birthdays of grown-up people are subject to prescribed
        rules," lady Feng expostulated; "but her present birthday is neither one of
        an adult nor that of an infant, and that's why I would like to deliberate
        with you!"}
}
{
}

\Verse{3}
{
    \zA{说着，一径去了，不在话 下。 且说湘云住了两日，便要回去，贾母因说：“等过了你宝姐姐的生日，看了戏， 再回去。”
        湘云听了，只得住下，又一面遣人回去，将自己旧日作的两件针线活计 取来，为宝钗生辰之仪。
        谁想贾母自见宝钗来了，喜他稳重和平，正值他才过第一个生辰，便自己捐资 二十两，唤了凤姐来，交与他备酒戏。 凤姐凑趣，笑道：“一个老祖宗，给孩子们
        作生日，不拘怎么着，谁还敢争？ 又办什么酒席呢？}
}
{
    \eA{Chia Lien upon hearing this remark, lowered his head and gave himself to
        protracted reflection. "You're indeed grown dull!" he cried; "why you've a
        precedent ready at hand to suit your case! Cousin Lin's birthday affords a
        precedent, and what you did in former years for cousin Lin, you can in this
        instance likewise do for cousin Hsüeh, and it will be all right." At these
        words lady Feng gave a sarcastic smile.}
}
{
}

\Verse{4}
{
    \zA{既高兴，要热闹，就说不得自己 花费几两老库里的体己。 这早晚找出这霉烂的二十两银子来做东，意思还叫我们赔
        上!果然拿不出来也罢了，金的银的圆的扁的压塌了箱子底，只是累？ 我们。 老祖 宗看看，谁不是你老人家的儿女？ 难道将来只有宝兄弟顶你老人家上五台山不成？
        那 些东西只留给他!我们虽不配使，也别太苦了我们。 这个够酒的够戏的呢？” 说的 满屋里都笑起来。
        贾母亦笑道：“你们听听这嘴!我也算会说的了，怎么说不过这 猴儿？}
}
{
    \eA{"Do you, pray, mean to insinuate," she added, "that I'm not aware of even
        this! I too had previously come, after some thought, to this conclusion;}
}
{
}

\Verse{5}
{
    \zA{你婆婆也不敢强嘴，你就和我？ 啊？ 的！” 凤姐笑道：“我婆婆也是一样的 疼宝玉，我也没处诉冤!倒说我强嘴！” 说着，又引贾母笑了一会。 贾母十分喜悦。
        到晚上，众人都在贾母前，定省之馀，大家娘儿们说笑时，贾母因问宝钗爱听何戏， 爱吃何物。 宝钗深知贾母年老之人，喜热闹戏文，爱吃甜烂之物，便总依贾母素喜
        者说了一遍。 贾母更加喜欢。 次日，先送过衣服玩物去，王夫人、凤姐、黛玉等诸 人皆有随分的，不须细说。}
}
{
    \eA{but old lady Chia explained, in my hearing yesterday, that having made
        inquiries about all their ages and their birthdays, she learnt that cousin
        Hsüeh would this year be fifteen, and that though this was not the birthday,
        which made her of age, she could anyhow well be regarded as being on the
        dawn of the year, in which she would gather up her hair, so that our dowager
        lady enjoined that her anniversary should, as a matter of course, be
        celebrated, unlike that of cousin Lin."}
}
{
}

\Verse{6}
{
    \zA{至二十一日，贾母内院搭了家常小巧戏台，定了一班新 出的小戏，昆弋两腔俱有。 就在贾母上房摆了几席家宴酒席，并无一个外客，只有
        薛姨妈、史湘云、宝钗是客，馀者皆是自己人。 这日早起，宝玉因不见黛玉，便到 他房中来寻，只见黛玉歪在炕上。 宝玉笑道：“起来吃饭去。 就开戏了，你爱听那
        一出？ 我好点。” 黛玉冷笑道：“你既这么说，你就特叫一班戏，拣我爱的唱给我 听，这会子犯不上借着光儿问我。”}
}
{
    \eA{"Well, in that case," Chia Lien suggested, "you had better make a few
        additions to what was done for cousin Lin!" "That's what I too am thinking
        of," lady Feng replied, "and that's why I'm asking your views; for were I,
        on my own hook, to add anything you would again feel hurt for my not have
        explained things to you." "That will do, that will do!" Chia Lien rejoined
        laughing, "none of these sham attentions for me!}
}
{
}

\Verse{7}
{
    \zA{宝玉笑道：“这有什么难的，明儿就叫一班子， 也叫他们借着咱们的光儿。” 一面说，一面拉他起来，携手出去。
        吃了饭，点戏时，贾母一面先叫宝钗点，宝钗推让一遍，无法，只得点了一出 《西游记》。 贾母自是喜欢。 又让薛姨妈，薛姨妈见宝钗点了，不肯再点。 贾母便
        特命凤姐点。 凤姐虽有邢王二夫人在前，但因贾母之命，不敢违拗，且知贾母喜热 闹更喜谑笑科诨，便先点了一出，却是《刘二当衣》。 贾母果真更又喜欢。}
}
{
    \eA{So long as you don't pry into my doings it will be enough; and will I go so
        far as to bear you a grudge?" With these words still in his mouth, he
        forthwith went off. But leaving him alone we shall now return to Shih
        Hsiang-yün. After a stay of a couple of days, her intention was to go back,
        but dowager lady Chia said: "Wait until after you have seen the theatrical
        performance, when you can return home."}
}
{
}

\Verse{8}
{
    \zA{然后便 命黛玉点，黛玉又让王夫人等先点。 贾母道：“今儿原是我特带着你们取乐，咱们 只管咱们的，别理他们。 我巴巴儿的唱戏摆酒，为他们呢？
        他们白听戏白吃已经便 宜了，还让他们点戏呢！” 说着，大家都笑。 黛玉方点了一出。 然后宝玉、史湘云、 迎、探、惜、李纨等俱各点了，按出扮演。
        至上酒席时，贾母又命宝钗点，宝钗点了一出《山门》。 宝玉道：“你只好点 这些戏。” 宝钗道：“你白听了这几年戏，那里知道这出戏，排场词藻都好呢。”}
}
{
    \eA{At this proposal, Shih Hsiang-yün felt constrained to remain, but she, at
        the same time, despatched a servant to her home to fetch two pieces of
        needlework, which she had in former days worked with her own hands, for a
        birthday present for Pao-ch'ai.}
}
{
}

\Verse{9}
{
    \zA{宝玉道：“我从来怕这些热闹戏。” 宝钗笑道：“要说这一出‘热闹’，你更不知 戏了。 你过来，我告诉你，这一出戏是一套《北点绛唇》，铿锵顿挫，那音律不用
        说是好了，那词藻中有只《寄生草》，极妙，你何曾知道！” 宝玉见说的这般好， 便凑近来央告：“好姐姐，念给我听听。” 宝钗便念给他听道： 漫？
        英雄泪，相离处士家。 谢慈悲剃度在莲台下。 没缘法转眼分离乍。 赤条条 来去无牵挂。 那里讨烟蓑雨笠卷单行？}
}
{
    \eA{Contrary to all expectations old lady Chia had, since the arrival of
        Pao-ch'ai, taken quite a fancy to her, for her sedateness and good nature,
        and as this happened to be the first birthday which she was about to
        celebrate (in the family) she herself readily contributed twenty taels
        which, after sending for lady Feng, she handed over to her, to make
        arrangements for a banquet and performance.}
}
{
}

\Verse{10}
{
    \zA{一任俺芒鞋破钵随缘化! 宝玉听了，喜的拍膝摇头，称赏不已； 又赞宝钗无书不知。 黛玉把嘴一撇道：“安 静些看戏吧!还没唱《山门》，你就《妆疯》了。”
        说的湘云也笑了。 于是大家看 戏，到晚方散。 贾母深爱那做小旦的和那做小丑的，因命人带进来，细看时，益发可怜见的。
        因问他年纪，那小旦才十一岁，小丑才九岁，大家叹息了一回。 贾母令人另拿些肉 果给他两个，又另赏钱。
        凤姐笑道：“这个孩子扮上活像一个人，你们再瞧不出来。”}
}
{
    \eA{"A venerable senior like yourself," lady Feng thereupon smiled and ventured,
        with a view to enhancing her good cheer, "is at liberty to celebrate the
        birthday of a child in any way agreeable to you, without any one presuming
        to raise any objection; but what's the use again of giving a banquet?}
}
{
}

\Verse{11}
{
    \zA{宝钗心内也知道，却点头不说； 宝玉也点了点头儿不敢说。 湘云便接口道：“我知 道，是像林姐姐的模样儿。” 宝玉听了，忙把湘云瞅了一眼。
        众人听了这话，留神 细看，都笑起来了，说：“果然像他！” 一时散了。 晚间，湘云便命翠缕把衣包收拾了。 翠缕道：“忙什么？ 等去的时候包也不迟。”
        湘云道：“明早就走，还在这里做什么？ ――看人家的脸子！” 宝玉听了这话，忙 近前说道：“好妹妹，你错怪了我。 林妹妹是个多心的人。}
}
{
    \eA{But since it be your good pleasure and your purpose to have it celebrated
        with éclat, you could, needless to say, your own self have spent several
        taels from the private funds in that old treasury of yours! But you now
        produce those twenty taels, spoiled by damp and mould, to play the hostess
        with, with the view indeed of compelling us to supply what's wanted!}
}
{
}

\Verse{12}
{
    \zA{别人分明知道，不肯说 出来，也皆因怕他恼。 谁知你不防头就说出来了，他岂不恼呢？ 我怕你得罪了人， 所以才使眼色。 你这会子恼了我，岂不辜负了我？
        要是别人，那怕他得罪了人，与 我何干呢？” 湘云摔手道：“你那花言巧语别望着我说。 我原不及你林妹妹。 别人 拿他取笑儿都使得，我说了就有不是。
        我本也不配和他说话：他是主子姑娘，我是 奴才丫头么。” 宝玉急的说道：“我倒是为你为出不是来了。 我要有坏心，立刻化 成灰，教万人拿脚踹！”}
}
{
    \eA{But hadn't you really been able to contribute any more, no one would have a
        word to say; but the gold and silver, round as well as flat, have with their
        heavy weight pressed down the bottom of the box! and your sole object is to
        harass us and to extort from us. But raise your eyes and look about you; who
        isn't your venerable ladyship's son and daughter?}
}
{
}

\Verse{13}
{
    \zA{湘云道：“大正月里，少信着嘴胡说这些没要紧的歪话! 你要说，你说给那些小性儿、行动爱恼人、会辖治你的人听去，别叫我啐你。” 说
        着，进贾母里间屋里，气忿忿的躺着去了。 宝玉没趣，只得又来找黛玉。 谁知才进门，便被黛玉推出来了，将门关上。 宝
        玉又不解何故，在窗外只是低声叫好妹妹好妹妹，黛玉总不理他。 宝玉闷闷的垂头 不语。 紫鹃却知端底，当此时料不能劝。 那宝玉只呆呆的站着。
        黛玉只当他回去了， 却开了门，只见宝玉还站在那里。}
}
{
    \eA{and is it likely, pray, that in the future there will only be cousin Pao-yü
        to carry you, our old lady, on his head, up the Wu T'ai Shan? You may keep
        all these things for him alone! but though we mayn't at present, deserve
        that anything should be spent upon us, you shouldn't go so far as to place
        us in any perplexities (by compelling us to subscribe).}
}
{
}

\Verse{14}
{
    \zA{黛玉不好再闭门，宝玉因跟进来，问道：“凡事 都有个原故，说出来人也不委屈。 好好的就恼，到底为什么起呢？” 黛玉冷笑道： “问我呢!我也不知为什么。
        我原是给你们取笑儿的，――拿着我比戏子，给众人 取笑儿！” 宝玉道：“我并没有比你，也并没有笑你，为什么恼我呢？” 黛玉道： “你还要比，你还要笑？
        你不比不笑，比人家比了笑了的还利害呢！” 宝玉听说， 无可分辩。 黛玉又道：“这还可恕。 你为什么又和云儿使眼色儿？ 这安的是什么心？}
}
{
    \eA{And is this now enough for wines, and enough for the theatricals?" As she
        bandied these words, every one in the whole room burst out laughing, and
        even dowager lady Chia broke out in laughter while she observed: "Do you
        listen to that mouth? I myself am looked upon as having the gift of the gab,
        but why is it that I can't talk in such a wise as to put down this monkey?}
}
{
}

\Verse{15}
{
    \zA{莫不是他和我玩，他就自轻自贱了？ 他是公侯的小姐，我原是民间的丫头。 他和我 玩，设如我回了口，那不是他自惹轻贱？ 你是这个主意不是？
        你却也是好心，只是那 一个不领你的情，一般也恼了。 你又拿我作情，倒说我‘小性儿、行动肯恼人’。
        你又怕他得罪了我，――我恼他与你何干，他得罪了我又与你何干呢？” 宝玉听了，方知才和湘云私谈，他也听见了。}
}
{
    \eA{Your mother-in-law herself doesn't dare to be so overbearing in her speech;
        and here you are jabber, jabber with me!" "My mother-in-law," explained lady
        Feng, "is also as fond of Pao-yü as you are, so much so that I haven't
        anywhere I could go and give vent to my grievances; and instead of (showing
        me some regard) you say that I'm overbearing in my speech!"}
}
{
}

\Verse{16}
{
    \zA{细想自己原为怕他二人恼了，故 在中间调停，不料自己反落了两处的数落，正合着前日所看《南华经》内“巧者劳
        而智者忧，无能者无所求，蔬食而遨游，泛若不系之舟”，又曰“山木自寇，源泉 自盗”等句，因此越想越无趣。 再细想来：“如今不过这几个人，尚不能应酬妥协，
        将来犹欲何为？” 想到其间，也不分辩，自己转身回房。 黛玉见他去了，便知回思 无趣，赌气去的，一言也不发，不禁自己越添了气，便说：“这一去，一辈子也别
        来了，也别说话！”}
}
{
    \eA{With these words, she again enticed dowager lady Chia to laugh for a while.
        The old lady continued in the highest of spirits, and, when evening came,
        and they all appeared in her presence to pay their obeisance, her ladyship
        made it a point, while the whole company of ladies and young ladies were
        engaged in chatting, to ascertain of Pao-ch'ai what play she liked to hear,
        and what things she fancied to eat.}
}
{
}

\Verse{17}
{
    \zA{那宝玉不理，竟回来，躺在床上，只是闷闷的。 袭人虽深知原 委，不敢就说，只得以别事来解说，因笑道：“今儿听了戏，又勾出几天戏来。 宝
        姑娘一定要还席的。” 宝玉冷笑道：“他还不还，与我什么相干？” 袭人见这话不 似往日，因又笑道：“这是怎么说呢？ 好好儿的大正月里，娘儿们姐儿们都喜喜欢
        欢的，你又怎么这个样儿了？” 宝玉冷笑道：“他们娘儿们姐儿们喜欢不喜欢，也 与我无干。” 袭人笑道：“大家随和儿，你也随点和儿不好？”}
}
{
    \eA{Pao-ch'ai was well aware that dowager lady Chia, well up in years though she
        was, delighted in sensational performances, and was partial to sweet and
        tender viands, so that she readily deferred, in every respect, to those
        things, which were to the taste of her ladyship, and enumerated a whole
        number of them, which made the old lady become the more exuberant.}
}
{
}

\Verse{18}
{
    \zA{宝玉道：“什么‘大 家彼此’？ 他们有‘大家彼此’，我只是赤条条无牵挂的！” 说到这句，不觉泪下。 袭人见这景况，不敢再说。
        宝玉细想这一句意味，不禁大哭起来。 翻身站起来，至 案边，提笔立占一偈云： 你证我证，心证意证。 是无有证，斯可云证。 无可云证，是立足境。
        写毕，自己虽解悟，又恐人看了不解，因又填一只《寄生草》，写在偈后。 又念了 一遍，自觉心中无有挂碍，便上床睡了。
        谁知黛玉见宝玉此番果断而去，假以寻袭人为由，来看动静。}
}
{
    \eA{And the next day, she was the first to send over clothes, nicknacks and such
        presents, while madame Wang and lady Feng, Tai-yü and the other girls, as
        well as the whole number of inmates had all presents for her, regulated by
        their degree of relationship, to which we need not allude in detail.}
}
{
}

\Verse{19}
{
    \zA{袭人回道：“已 经睡了。” 黛玉听了，就欲回去，袭人笑道：“姑娘请站着，有一个字帖儿，瞧瞧 写的是什么话。” 便将宝玉方才所写的拿给黛玉看。
        黛玉看了，知是宝玉为一时感 忿而作，不觉又可笑又可叹。 便向袭人道：“作的是个玩意儿，无甚关系的。” 说 毕，便拿了回房去。 次日，和宝钗湘云同看。
        宝钗念其词曰： 无我原非你，从他不解伊。 肆行无碍凭来去。 茫茫着甚悲愁喜，纷纷说甚亲疏 密。 从前碌碌却因何？}
}
{
    \eA{When the 21st arrived, a stage of an ordinary kind, small but yet handy, was
        improvised in dowager lady Chia's inner court, and a troupe of young actors,
        who had newly made their début, was retained for the nonce, among whom were
        both those who could sing tunes, slow as well as fast.}
}
{
}

\Verse{20}
{
    \zA{到如今回头试想真无趣! 看毕，又看那偈语，因笑道：“这是我的不是了。 我昨儿一支曲子，把他这个话惹 出来。
        这些道书机锋，最能移性的，明儿认真说起这些疯话，存了这个念头，岂不 是从我这支曲子起的呢？ 我成了个罪魁了！” 说着，便撕了个粉碎，递给丫头们，
        叫快烧了。 黛玉笑道：“不该撕了，等我问他，你们跟我来，包管叫他收了这个痴 心。” 三人说着，过来见了宝玉。
        黛玉先笑道：“宝玉，我问你：至贵者宝，至坚者 玉。 尔有何贵？ 尔有何坚？”}
}
{
    \eA{In the drawing rooms of the old lady were then laid out several tables for a
        family banquet and entertainment, at which there was not a single outside
        guest; and with the exception of Mrs. Hsüeh, Shih Hsiang-yün, and Pao-ch'ai,
        who were visitors, the rest were all inmates of her household.}
}
{
}

\Verse{21}
{
    \zA{宝玉竟不能答。 二人笑道：“这样愚钝，还参禅呢！” 湘云也拍手笑道：“宝哥哥可输了。” 黛玉又道：“你道‘无可云证，是立足境’，
        固然好了，只是据我看来，还未尽善。 我还续两句云：‘无立足境，方是干净。’” 宝钗道：“实在这方悟彻。 当日南宗六祖惠能初寻师至韶州，闻五祖弘忍在黄梅，
        他便充作火头僧。 五祖欲求法嗣，令诸僧各出一偈，上座神秀说道：‘身是菩提树， 心如明镜台。 时时勤拂拭，莫使有尘埃。’}
}
{
    \eA{On this day, Pao-yü failed, at any early hour, to see anything of Lin
        Tai-yü, and coming at once to her rooms in search of her, he discovered her
        reclining on the stove-couch. "Get up," Pao-yü pressed her with a smile,
        "and come and have breakfast, for the plays will commence shortly; but
        whichever plays you would like to listen to, do tell me so that I may be
        able to choose them." Tai-yü smiled sarcastically.}
}
{
}

\Verse{22}
{
    \zA{惠能在厨房舂米，听了道：‘美则美矣， 了则未了。’ 因自念一偈曰：‘菩提本非树，明镜亦非台。 本来无一物，何处染尘 埃？’ 五祖便将衣钵传给了他。
        今儿这偈语亦同此意了。 只是方才这句机锋，尚未 完全了结，这便丢开手不成？” 黛玉笑道：“他不能答就算输了，这会子答上了也 不为出奇了。
        只是以后再不许谈禅了。 连我们两个人所知所能的，你还不知不能呢， 还去参什么禅呢！” 宝玉自己以为觉悟，不想忽被黛玉一问，便不能答；}
}
{
    \eA{"In that case," she rejoined, "you had better specially engage a troupe and
        select those I like sung for my benefit; for on this occasion you can't be
        so impertinent as to make use of their expense to ask me what I like!"
        "What's there impossible about this?" Pao-yü answered smiling; "well,
        to-morrow I'll readily do as you wish, and ask them too to make use of what
        is yours and mine."}
}
{
}

\Verse{23}
{
    \zA{宝钗又比 出语录来，此皆素不见他们所能的。 自己想了一想：“原来他们比我的知觉在先， 尚未解悟，我如今何必自寻苦恼。”
        想毕，便笑道：“谁又参禅，不过是一时的玩 话儿罢了。” 说罢，四人仍复如旧。 忽然人报娘娘差人送出一个灯谜来，命他们大家去猜，猜后每人也作一个送进 去。
        四人听说，忙出来至贾母上房，只见一个小太监，拿了一盏四角平头白纱灯， 专为灯谜而制，上面已有了一个，众人都争看乱猜。}
}
{
    \eA{As he passed this remark, he pulled her up, and taking her hand in his own,
        they walked out of the room and came and had breakfast. When the time
        arrived to make a selection of the plays, dowager lady Chia of her own
        motion first asked Pao-ch'ai to mark off those she liked;}
}
{
}

\Verse{24}
{
    \zA{小太监又下谕道：“众小姐猜 着，不要说出来，每人只暗暗的写了，一齐封送进去，候娘娘自验是否。” 宝钗听
        了，近前一看，是一首七言绝句，并无新奇，口中少不得称赞，只说“难猜”，故 意寻思。 其实一见早猜着了。 宝玉、黛玉、湘云、探春四个人也都解了，各自暗暗
        的写了。 一并将贾环贾兰等传来，一齐各揣心机猜了，写在纸上，然后各人拈一物 作成一谜，恭楷写了，挂于灯上。}
}
{
    \eA{and though for a time Pao-ch'ai declined, yielding the choice to others, she
        had no alternative but to decide, fixing upon a play called, "the Record of
        the Western Tour," a play of which the old lady was herself very fond.}
}
{
}

\Verse{25}
{
    \zA{太监去了，至晚出来，传谕道：“前日娘娘所制，俱已猜着，惟二小姐与三爷 猜的不是。 小姐们作的也都猜了，不知是否？” 说着，也将写的拿出来，也有猜着
        的，也有猜不着的。 太监又将颁赐之物送与猜着之人，每人一个宫制诗筒，一柄茶 筅，独迎春贾环二人未得。 迎春自以为玩笑小事，并不介意； 贾环便觉得没趣。
        且 又听太监说：“三爷所作这个不通，娘娘也没猜，叫我带回问三爷是个什么。”}
}
{
    \eA{Next in order, she bade lady Feng choose, and lady Feng, had, after all, in
        spite of madame Wang ranking before her in precedence, to consider old lady
        Chia's request, and not to presume to show obstinacy by any disobedience.}
}
{
}

\Verse{26}
{
    \zA{众 人听了，都来看他作的是什么，――写道： 大哥有角只八个，二哥有角只两根。 大哥只在床上坐，二哥爱在房上蹲。 众人看了，大发一笑。
        贾环只得告诉太监说：“是一个枕头，一个兽头。” 太监记 了，领茶而去。 贾母见元春这般有兴，自己一发喜乐，便命速作一架小巧精致围屏灯来，设于
        堂屋，命他姊妹们各自暗暗的做了，写出来粘在屏上； 然后预备下香茶细果以及各 色玩物，为猜着之贺。 贾政朝罢，见贾母高兴，况在节间，晚上也来承欢取乐。}
}
{
    \eA{But as she knew well enough that her ladyship had a penchant for what was
        exciting, and that she was still more partial to jests, jokes, epigrams, and
        buffoonery, she therefore hastened to precede (madame Wang) and to choose a
        play, which was in fact no other than "Liu Erh pawns his clothes." Dowager
        lady Chia was, of course, still more elated. And after this she speedily
        went on to ask Tai-yü to choose.}
}
{
}

\Verse{27}
{
    \zA{上 面贾母、贾政、宝玉一席； 王夫人、宝钗、黛玉、湘云又一席，迎春、探春、惜春 三人又一席，俱在下面。 地下老婆丫鬟站满。
        李宫裁王熙凤二人在里间又一席。 贾 政因不见贾兰，便问：“怎么不见兰哥儿？” 地下女人们忙进里间问李氏，李氏起
        身笑着回道：“他说方才老爷并没叫他去，他不肯来。” 女人们回复了贾政，众人 都笑说：“天生的牛心拐孤！” 贾政忙遣贾环和个女人将贾兰唤来，贾母命他在身
        边坐了，抓果子给他吃，大家说笑取乐。}
}
{
    \eA{Tai-yü likewise concedingly yielded her turn in favour of madame Wang and
        the other seniors, to make their selections before her, but the old lady
        expostulated. "To-day," she said, "is primarily an occasion, on which I've
        brought all of you here for your special recreation; and we had better look
        after our own selves and not heed them!}
}
{
}

\Verse{28}
{
    \zA{往常间只有宝玉长谈阔论，今日贾政在这 里，便唯唯而已。 馀者，湘云虽系闺阁弱质，却素喜谈论，今日贾政在席，也自？ 口禁语； 黛玉本性娇懒，不肯多话；
        宝钗原不妄言轻动，便此时亦是坦然自若：故 此一席，虽是家常取乐，反见拘束。 贾母亦知因贾政一人在此所致，酒过三巡，便撵贾政去歇息。 贾政亦知贾母之
        意，撵了他去好让他姊妹兄弟们取乐，因陪笑道：“今日原听见老太太这里大设春 灯雅谜，故也备了彩礼酒席，特来入会。}
}
{
    \eA{For have I, do you imagine, gone to the trouble of having a performance and
        laying a feast for their special benefit? they're already reaping benefit
        enough by being in here, listening to the plays and partaking of the
        banquet, when they have no right to either; and are they to be pressed
        further to make a choice of plays?" At these words, the whole company had a
        hearty laugh;}
}
{
}

\Verse{29}
{
    \zA{何疼孙子孙女之心，便不略赐与儿子半 点？” 贾母笑道：“你在这里，他们都不敢说笑，没的倒叫我闷的慌。 你要猜谜儿， 我说一个你猜，猜不着是要罚的。”
        贾政忙笑道：“自然受罚。 若猜着了，也要领 赏呢。” 贾母道：“这个自然。” 便念道： 猴子身轻站树梢。 打一果名。
        贾政已知是荔枝，故意乱猜，罚了许多东西，然后方猜着了，也得了贾母的东西。 然后也念一个灯谜与贾母猜。 念道： 身自端方，体自坚硬。 虽不能言，有言必应。
        打一用物。}
}
{
    \eA{after which, Tai-yü, at length, marked off a play; next in order following
        Pao-yü, Shih Hsiang-yün, Ying-ch'un, T'an Ch'un, Hsi Ch'un, widow Li Wan,
        and the rest, each and all of whom made a choice of plays, which were sung
        in the costumes necessary for each.}
}
{
}

\Verse{30}
{
    \zA{说毕，便悄悄的说与宝玉，宝玉会意，又悄悄的告诉了贾母。 贾母想了一想，果然 不差，便说：“是砚台。” 贾政笑道：“到底是老太太，一猜就是。”
        回头说：“快 把贺彩献上来。” 地下妇女答应一声，大盘小盒，一齐捧上。 贾母逐件看去，都是 灯节下所用所玩新巧之物，心中甚喜，遂命：“给你老爷斟酒。”
        宝玉执壶，迎春 送酒。 贾母因说：“你瞧瞧那屏上，都是他姐儿们做的，再猜一猜我听。”}
}
{
    \eA{When the time came to take their places at the banquet, dowager lady Chia
        bade Pao-ch'ai make another selection, and Pao-ch'ai cast her choice upon
        the play: "Lu Chih-shen, in a fit of drunkenness stirs up a disturbance up
        the Wu T'ai mountain;" whereupon Pao-yü interposed, with the remark: "All
        you fancy is to choose plays of this kind;"}
}
{
}

\Verse{31}
{
    \zA{贾政答应，起身走至屏前，只见第一个是元妃的，写着道： 能使妖魔胆尽摧，身如束帛气如雷。 一声震得人方恐，回首相看已化灰。 打一玩物。
        贾政道：“这是爆竹吗？” 宝玉答道：“是。” 贾政又看迎春的，道： 天运人功理不穷，有功无运也难逢。 因何镇日纷纷乱？ 只为阴阳数不通。 打一用物。
        贾政道：“是算盘？” 迎春笑道：“是。” 又往下看，是探春的，道： 阶下儿童仰面时，清明妆点最堪宜。 游丝一断浑无力，莫向东风怨别离。 打一玩物。}
}
{
    \eA{to which Pao-ch'ai rejoined, "You've listened to plays all these years to no
        avail! How could you know the beauties of this play? the stage effect is
        grand, but what is still better are the apt and elegant passages in it."
        "I've always had a dread of such sensational plays as these!" Pao-yü
        retorted.}
}
{
}

\Verse{32}
{
    \zA{贾政道：“好像风筝。” 探春道：“是。” 贾政再往下看，是黛玉的，道： 朝罢谁携两袖烟？ 琴边衾里两无缘。 晓筹不用鸡人报，五夜无烦侍女添。
        焦首朝朝还暮暮，煎心日日复年年。 光阴荏苒须当惜，风雨阴晴任变迁。 打一用物。 贾政道：“这个莫非是更香？” 宝玉代言道：“是。” 贾政又看道：
        南面而坐，北面而朝。 “象忧亦忧，象喜亦喜。” 打一用物。 贾政道：“好，好!如猜镜子，妙极！” 宝玉笑回道：“是。”}
}
{
    \eA{"If you call this play sensational," Pao-ch'ai smilingly expostulated, "well
        then you may fitly be looked upon as being no connoisseur of plays. But come
        over and I'll tell you. This play constitutes one of a set of books,
        entitled the 'Pei Tien Peng Ch'un,' which, as far as harmony, musical rests
        and closes, and tune go, is, it goes without saying, perfect;}
}
{
}

\Verse{33}
{
    \zA{贾政道：“这一个 却无名字，是谁做的？” 贾母道：“这个大约是宝玉做的？” 贾政就不言语。 往下 再看宝钗的，道是： 有眼无珠腹内空，荷花出水喜相逢。
        梧桐叶落分离别，恩爱夫妻不到冬。 打一用物。 贾政看完，心内自忖道：“此物还倒有限，只是小小年纪，作此等言语，更觉不祥。 看来皆非福寿之辈。”
        想到此处，甚觉烦闷，大有悲戚之状，只是垂头沉思。}
}
{
    \eA{but there's among the elegant compositions a ballad entitled: 'the Parasitic
        Plant,' written in a most excellent style; but how could you know anything
        about it?" Pao-yü, upon hearing her speak of such points of beauty, hastily
        drew near to her. "My dear cousin," he entreated, "recite it and let me hear
        it!"}
}
{
}

\Verse{34}
{
    \zA{贾母 见贾政如此光景，想到他身体劳乏，又恐拘束了他众姊妹，不得高兴玩耍，便对贾 政道：“你竟不必在这里了，歇着去罢。 让我们再坐一会子，也就散了。”
        贾政一 闻此言，连忙答应几个“是”，又勉强劝了贾母一回酒，方才退出去了。 回至房中， 只是思索，翻来覆去，甚觉凄惋。
        这里贾母见贾政去了，便道：“你们乐一乐罢。” 一语未了，只见宝玉跑至围 屏灯前，指手画脚，信口批评：“这个这一句不好。” “那个破的不恰当。” 如同
        开了锁的猴儿一般。}
}
{
    \eA{Whereupon Pao-ch'ai went on as follows:  My manly tears I will not wipe
        away,  But from this place, the scholar's home, I'll stray. The bonze for
        mercy I shall thank; under the lotus altar shave my  pate; With Yüan to be
        the luck I lack; soon in a twinkle we shall separate,  And needy and forlorn
        I'll come and go, with none to care about my  fate. Thither shall I a
        suppliant be for a fog wrapper and rain hat;}
}
{
}

\Verse{35}
{
    \zA{黛玉便道：“还像方才大家坐着，说说笑笑，岂不斯文些儿？” 凤姐儿自里间屋里出来，插口说道：“你这个人，就该老爷每日合你寸步儿不离才 好。
        刚才我忘了，为什么不当着老爷，撺掇着叫你作诗谜儿？ 这会子不怕你不出汗 呢。” 说的宝玉急了，扯着凤姐儿厮缠了一会。 贾母又和李宫裁并众姊妹等说笑了
        一会子，也觉有些困倦，听了听，已交四鼓了。 因命将食物撤去，赏给众人，遂起 身道：“我们歇着罢。 明日还是节呢，该当早些起来。 明日晚上再玩罢。”}
}
{
    \eA{my  warrant I shall roll,  And listless with straw shoes and broken bowl,
        wherever to convert my  fate may be, I'll stroll. As soon as Pao-yü had
        listened to her recital, he was so full of enthusiasm, that, clapping his
        knees with his hands, and shaking his head, he gave vent to incessant
        praise; after which he went on to extol Pao-ch'ai, saying: "There's no book
        that you don't know."}
}
{
}

\Verse{36}
{
    \zA{于是众 人方慢慢的散去。 未知次日如何，且听下回分解。 (本章完)}
}
{
    \eA{"Be quiet, and listen to the play," Lin Tai-yü urged; "they haven't yet sung
        about the mountain gate, and you already pretend to be mad!" At these words,
        Hsiang-yün also laughed. But, in due course, the whole party watched the
        performance until evening, when they broke up.}
}
{
}

\Verse{37}
{
    \zA{}
}
{
    \eA{Dowager lady Chia was so very much taken with the young actor, who played
        the role of a lady, as well as with the one who acted the buffoon, that she
        gave orders that they should be brought in; and, as she looked at them
        closely, she felt so much the more interest in them, that she went on to
        inquire what their ages were. And when the would-be lady (replied) that he
        was just eleven, while the would-be buffoon (explained) that he was just
        nine, the whole company gave vent for a time to expressions of sympathy with
        their lot; while dowager lady Chia bade servants bring a fresh supply of
        meats and fruits for both of them, and also gave them, besides their wages,
        two tiaos as a present.}
}
{
}

\Verse{38}
{
    \zA{}
}
{
    \eA{"This lad," lady Feng observed smiling, "is when dressed up (as a girl), a
        living likeness of a certain person; did you notice it just now?" Pao-ch'ai
        was also aware of the fact, but she simply nodded her head assentingly and
        did not say who it was. Pao-yü likewise expressed his assent by shaking his
        head, but he too did not presume to speak out. Shih Hsiang-yün, however,
        readily took up the conversation. "He resembles," she interposed, "cousin
        Lin's face!" When this remark reached Pao-yü's ear, he hastened to cast an
        angry scowl at Hsiang-yün, and to make her a sign; while the whole party,
        upon hearing what had been said, indulged in careful and minute scrutiny of
        (the lad); and as they all began to laugh: "The resemblance is indeed
        striking!" they exclaimed. After a while, they parted;}
}
{
}

\Verse{39}
{
    \zA{}
}
{
    \eA{and when evening came Hsiang-yün directed Ts'ui Lü to pack up her clothes.
        "What's the hurry?" Ts'ui Lü asked. "There will be ample time to pack up, on
        the day on which we go!" "We'll go to-morrow," Hsiang-yün rejoined; "for
        what's the use of remaining here any longer--to look at people's mouths and
        faces?" Pao-yü, at these words, lost no time in pressing forward. "My dear
        cousin," he urged; "you're wrong in bearing me a grudge! My cousin Lin is a
        girl so very touchy, that though every one else distinctly knew (of the
        resemblance), they wouldn't speak out; and all because they were afraid that
        she would get angry; but unexpectedly out you came with it, at a moment when
        off your guard; and how ever couldn't she but feel hurt?}
}
{
}

\Verse{40}
{
    \zA{}
}
{
    \eA{and it's because I was in dread that you would give offence to people that I
        then winked at you; and now here you are angry with me; but isn't that being
        ungrateful to me? Had it been any one else, would I have cared whether she
        had given offence to even ten; that would have been none of my business!"
        Hsiang-yün waved her hand: "Don't," she added, "come and tell me these
        flowery words and this specious talk, for I really can't come up to your
        cousin Lin. If others poke fun at her, they all do so with impunity, while
        if I say anything, I at once incur blame. The fact is I shouldn't have
        spoken of her, undeserving as I am;}
}
{
}

\Verse{41}
{
    \zA{}
}
{
    \eA{and as she's the daughter of a master, while I'm a slave, a mere servant
        girl, I've heaped insult upon her!" "And yet," pleaded Pao-yü, full of
        perplexity, "I had done it for your sake; and through this, I've come in for
        reproach. But if it were with an evil heart I did so, may I at once become
        ashes, and be trampled upon by ten thousands of people!" "In this felicitous
        firstmonth," Hsiang-yün remonstrated, "you shouldn't talk so much reckless
        nonsense! All these worthless despicable oaths, disjointed words, and
        corrupt language, go and tell for the benefit of those mean sort of people,
        who in everything take pleasure in irritating others, and who keep you under
        their thumb! But mind don't drive me to spit contemptuously at you."}
}
{
}

\Verse{42}
{
    \zA{}
}
{
    \eA{As she gave utterance to these words, she betook herself in the inner room
        of dowager lady Chia's suite of apartments, where she lay down in high
        dudgeon, and, as Pao-yü was so heavy at heart, he could not help coming
        again in search of Tai-yü; but strange to say, as soon as he put his foot
        inside the doorway, he was speedily hustled out of it by Tai-yü, who shut
        the door in his face. Pao-yü was once more unable to fathom her motives, and
        as he stood outside the window, he kept on calling out: "My dear cousin," in
        a low tone of voice; but Tai-yü paid not the slightest notice to him so that
        Pao-yü became so melancholy that he drooped his head, and was plunged in
        silence.}
}
{
}

\Verse{43}
{
    \zA{}
}
{
    \eA{And though Hsi Jen had, at an early hour, come to know the circumstances,
        she could not very well at this juncture tender any advice. Pao-yü remained
        standing in such a vacant mood that Tai-yü imagined that he had gone back;
        but when she came to open the door she caught sight of Pao-yü still waiting
        in there; and as Tai-yü did not feel justified to again close the door,
        Pao-yü consequently followed her in. "Every thing has," he observed, "a why
        and a wherefore; which, when spoken out, don't even give people pain; but
        you will rush into a rage, and all without any rhyme! but to what really
        does it owe its rise?" "It's well enough, after all, for you to ask me,"
        Tai-yü rejoined with an indifferent smile, "but I myself don't know why!}
}
{
}

\Verse{44}
{
    \zA{}
}
{
    \eA{But am I here to afford you people amusement that you will compare me to an
        actress, and make the whole lot have a laugh at me?" "I never did liken you
        to anything," Pao-yü protested, "neither did I ever laugh at you! and why
        then will you get angry with me?" "Was it necessary that you should have
        done so much as made the comparison," Tai-yü urged, "and was there any need
        of even any laughter from you? why, though you mayn't have likened me to
        anything, or had a laugh at my expense, you were, yea more dreadful than
        those who did compare me (to a singing girl) and ridiculed me!" Pao-yü could
        not find anything with which to refute the argument he had just heard, and
        Tai-yü went on to say. "This offence can, anyhow, be condoned; but, what is
        more, why did you also wink at Yün Erh?}
}
{
}

\Verse{45}
{
    \zA{}
}
{
    \eA{What was this idea which you had resolved in your mind? wasn't it perhaps
        that if she played with me, she would be demeaning herself, and making
        herself cheap? She's the daughter of a duke or a marquis, and we forsooth
        the mean progeny of a poor plebeian family; so that, had she diverted
        herself with me, wouldn't she have exposed herself to being depreciated, had
        I, perchance, said anything in retaliation? This was your idea wasn't it?
        But though your purpose was, to be sure, honest enough, that girl wouldn't,
        however, receive any favours from you, but got angry with you just as much
        as I did; and though she made me also a tool to do you a good turn, she, on
        the contrary, asserts that I'm mean by nature and take pleasure in
        irritating people in everything!}
}
{
}

\Verse{46}
{
    \zA{}
}
{
    \eA{and you again were afraid lest she should have hurt my feelings, but, had I
        had a row with her, what would that have been to you? and had she given me
        any offence, what concern would that too have been of yours?" When Pao-yü
        heard these words, he at once became alive to the fact that she too had lent
        an ear to the private conversation he had had a short while back with
        Hsiang-yün: "All because of my, fears," he carefully mused within himself,
        "lest these two should have a misunderstanding, I was induced to come
        between them, and act as a mediator;}
}
{
}

\Verse{47}
{
    \zA{}
}
{
    \eA{but I myself have, contrary to my hopes, incurred blame and abuse on both
        sides! This just accords with what I read the other day in the Nan Hua
        Ching. 'The ingenious toil, the wise are full of care; the good-for-nothing
        seek for nothing, they feed on vegetables, and roam where they list; they
        wander purposeless like a boat not made fast!' 'The mountain trees,' the
        text goes on to say, 'lead to their own devastation; the spring (conduces)
        to its own plunder; and so on." And the more he therefore indulged in
        reflection, the more depressed he felt. "Now there are only these few
        girls," he proceeded to ponder minutely, "and yet, I'm unable to treat them
        in such a way as to promote perfect harmony; and what will I forsooth do by
        and by (when there will be more to deal with)!"}
}
{
}

\Verse{48}
{
    \zA{}
}
{
    \eA{When he had reached this point in his cogitations, (he decided) that it was
        really of no avail to agree with her, so that turning round, he was making
        his way all alone into his apartments; but Lin Tai-yü, upon noticing that he
        had left her side, readily concluded that reflection had marred his spirits
        and that he had so thoroughly lost his temper as to be going without even
        giving vent to a single word, and she could not restrain herself from
        feeling inwardly more and more irritated. "After you've gone this time," she
        hastily exclaimed, "don't come again, even for a whole lifetime;}
}
{
}

\Verse{49}
{
    \zA{}
}
{
    \eA{and I won't have you either so much as speak to me!" Pao-yü paid no heed to
        her, but came back to his rooms, and laying himself down on his bed, he kept
        on muttering in a state of chagrin; and though Hsi Jen knew full well the
        reasons of his dejection, she found it difficult to summon up courage to say
        anything to him at the moment, and she had no alternative but to try and
        distract him by means of irrelevant matters. "The theatricals which you've
        seen to-day," she consequently observed smiling, "will again lead to
        performances for several days, and Miss Pao-ch'ai will, I'm sure, give a
        return feast."}
}
{
}

\Verse{50}
{
    \zA{}
}
{
    \eA{"Whether she gives a return feast or not," Pao-yü rejoined with an apathetic
        smirk, "is no concern of mine!" When Hsi Jen perceived the tone, so unlike
        that of other days, with which these words were pronounced: "What's this
        that you're saying?" she therefore remarked as she gave another smile. "In
        this pleasant and propitious first moon, when all the ladies and young
        ladies are in high glee, how is it that you're again in a mood of this
        sort?" "Whether the ladies and my cousins be in high spirits or not," Pao-yü
        replied forcing a grin, "is also perfectly immaterial to me."}
}
{
}

\Verse{51}
{
    \zA{}
}
{
    \eA{"They are all," Hsi Jen added, smilingly, "pleasant and agreeable, and were
        you also a little pleasant and agreeable, wouldn't it conduce to the
        enjoyment of the whole company?" "What about the whole company, and they and
        I?" Pao-yü urged. "They all have their mutual friendships; while I, poor
        fellow, all forlorn, have none to care a rap for me." His remarks had
        reached this clause, when inadvertently the tears trickled down; and Hsi Jen
        realising the state of mind he was in, did not venture to say anything
        further.}
}
{
}

\Verse{52}
{
    \zA{}
}
{
    \eA{But as soon as Pao-yü had reflected minutely over the sense and import of
        this sentence, he could not refrain from bursting forth into a loud fit of
        crying, and, turning himself round, he stood up, and, drawing near the
        table, he took up the pencil, and eagerly composed these enigmatical lines:
        If thou wert me to test, and I were thee to test,  Our hearts were we to
        test, and our minds to test,  When naught more there remains for us to test
        That will yea very well be called a test,  And when there's naught to put,
        we could say, to the test,  We will a place set up on which our feet to
        rest. After he had finished writing, he again gave way to fears that though
        he himself could unfold their meaning, others, who came to peruse these
        lines, would not be able to fathom them, and he also went on consequently to
        indite another stanza, in imitation of the "Parasitic Plant," which he
        inscribed at the close of the enigma;}
}
{
}

\Verse{53}
{
    \zA{}
}
{
    \eA{and when he had read it over a second time, he felt his heart so free of all
        concern that forthwith he got into his bed, and went to sleep. But, who
        would have thought it, Tai-yü, upon seeing Pao-yü take his departure in such
        an abrupt manner, designedly made use of the excuse that she was bent upon
        finding Hsi Jen, to come round and see what he was up to. "He's gone to
        sleep long ago!" Hsi Jen replied. At these words, Tai-yü felt inclined to
        betake herself back at once; but Hsi Jen smiled and said: "Please stop,
        miss. Here's a slip of paper, and see what there is on it!" and speedily
        taking what Pao-yü had written a short while back, she handed it over to
        Tai-yü to examine. Tai-yü, on perusal, discovered that Pao-yü had composed
        it, at the spur of the moment, when under the influence of resentment;}
}
{
}

\Verse{54}
{
    \zA{}
}
{
    \eA{and she could not help thinking it both a matter of ridicule as well as of
        regret; but she hastily explained to Hsi Jen: "This is written for fun, and
        there's nothing of any consequence in it!" and having concluded this remark,
        she readily took it along with her to her room, where she conned it over in
        company with Hsiang-yün; handing it also the next day to Pao-ch'ai to
        peruse. The burden of what Pao-ch'ai read was:  In what was no concern of
        mine, I should to thee have paid no heed,  For while I humour this, that one
        to please I don't succeed! Act as thy wish may be! go, come whene'er thou
        list; 'tis naught to  me. Sorrow or joy, without limit or bound, to indulge
        thou art free! What is this hazy notion about relatives distant or close?}
}
{
}

\Verse{55}
{
    \zA{}
}
{
    \eA{For what purpose have I for all these days racked my heart with woes? Even
        at this time when I look back and think, my mind no pleasure  knows. After
        having finished its perusal, she went on to glance at the Buddhistic stanza,
        and smiling: "This being," she soliloquised; "has awakened to a sense of
        perception; and all through my fault, for it's that ballad of mine yesterday
        which has incited this! But the subtle devices in all these rationalistic
        books have a most easy tendency to unsettle the natural disposition, and if
        to-morrow he does actually get up, and talk a lot of insane trash, won't his
        having fostered this idea owe its origin to that ballad of mine; and shan't
        I have become the prime of all guilty people?"}
}
{
}

\Verse{56}
{
    \zA{}
}
{
    \eA{Saying this, she promptly tore the paper, and, delivering the pieces to the
        servant girls, she bade them go at once and burn them. "You shouldn't have
        torn it!" Tai-yü remonstrated laughingly. "But wait and I'll ask him about
        it! so come along all of you, and I vouch I'll make him abandon that idiotic
        frame of mind and that depraved language." The three of them crossed over,
        in point of fact, into Pao-yü's room, and Tai-yü was the first to smile and
        observe. "Pao-yü, may I ask you something? What is most valuable is a
        precious thing; and what is most firm is jade, but what value do you possess
        and what firmness is innate in you?"}
}
{
}

\Verse{57}
{
    \zA{}
}
{
    \eA{But as Pao-yü could not, say anything by way of reply, two of them remarked
        sneeringly: "With all this doltish bluntness of his will he after all absorb
        himself in abstraction?" While Hsiang-yün also clapped her hands and
        laughed, "Cousin Pao has been discomfited." "The latter part of that
        apothegm of yours," Tai-yü continued, "says:  "We would then find some place
        on which our feet to rest. "Which is certainly good; but in my view, its
        excellence is not as yet complete! and I should still tag on two lines at
        its close;" as she proceeded to recite:  "If we do not set up some place on
        which our feet to rest,  For peace and freedom then it will be best." "There
        should, in very truth, be this adjunct to make it thoroughly explicit!"
        Pao-ch'ai added.}
}
{
}

\Verse{58}
{
    \zA{}
}
{
    \eA{"In days of yore, the sixth founder of the Southern sect, Hui Neng, came,
        when he went first in search of his patron, in the Shao Chou district; and
        upon hearing that the fifth founder, Hung Jen, was at Huang Mei, he readily
        entered his service in the capacity of Buddhist cook; and when the fifth
        founder, prompted by a wish to select a Buddhistic successor, bade his
        neophytes and all the bonzes to each compose an enigmatical stanza, the one
        who occupied the upper seat, Shen Hsiu, recited:  "A P'u T'i tree the body
        is, the heart so like a stand of mirror  bright,  On which must needs, by
        constant careful rubbing, not be left dust to  alight!}
}
{
}

\Verse{59}
{
    \zA{}
}
{
    \eA{"And Hui Neng, who was at this time in the cook-house pounding rice,
        overheard this enigma. 'Excellent, it is excellent,' he ventured, 'but as
        far as completeness goes it isn't complete;' and having bethought himself of
        an apothegm: 'The P'u T'i, (an expression for Buddha or intelligence),' he
        proceeded, 'is really no tree; and the resplendent mirror, (Buddhistic term
        for heart), is likewise no stand; and as, in fact, they do not constitute
        any tangible objects, how could they be contaminated by particles of dust?'
        Whereupon the fifth founder at once took his robe and clap-dish and handed
        them to him.}
}
{
}

\Verse{60}
{
    \zA{}
}
{
    \eA{Well, the text now of this enigma presents too this identical idea, for the
        simple fact is that those lines full of subtleties of a short while back are
        not, as yet, perfected or brought to an issue, and do you forsooth readily
        give up the task in this manner?" "He hasn't been able to make any reply,"
        Tai-yü rejoined sneeringly, "and must therefore be held to be discomfited;
        but were he even to make suitable answer now, there would be nothing out of
        the common about it! Anyhow, from this time forth you mustn't talk about
        Buddhistic spells, for what even we two know and are able to do, you don't
        as yet know and can't do; and do you go and concern yourself with
        abstraction?"}
}
{
}

\Verse{61}
{
    \zA{}
}
{
    \eA{Pao-yü had, in his own mind, been under the impression that he had attained
        perception, but when he was unawares and all of a sudden subjected to this
        question by Tai-yü, he soon found it beyond his power to give any ready
        answer. And when Pao-ch'ai furthermore came out with a religious
        disquisition, by way of illustration, and this on subjects, in all of which
        he had hitherto not seen them display any ability, he communed within
        himself: "If with their knowledge, which is indeed in advance of that of
        mine, they haven't, as yet, attained perception, what need is there for me
        now to bring upon myself labour and vexation?"}
}
{
}

\Verse{62}
{
    \zA{}
}
{
    \eA{"Who has, pray," he hastily inquired smilingly, after arriving at the end of
        his reflections, "indulged in Buddhistic mysteries? what I did amounts to
        nothing more than nonsensical trash, written, at the spur of the moment, and
        nothing else." At the close of this remark all four came to be again on the
        same terms as of old; but suddenly a servant announced that the Empress
        (Yüan Ch'un) had despatched a messenger to bring over a lantern-conundrum
        with the directions that they should all go and guess it, and that after
        they had found it out, they should each also devise one and send it in.}
}
{
}

\Verse{63}
{
    \zA{}
}
{
    \eA{At these words, the four of them left the room with hasty step, and
        adjourned into dowager lady Chia's drawing room, where they discovered a
        young eunuch, holding a four-cornered, flat-topped lantern, of white gauze,
        which had been specially fabricated for lantern riddles. On the front side,
        there was already a conundrum, and the whole company were vying with each
        other in looking at it and making wild guesses; when the young eunuch went
        on to transmit his orders, saying: "Young ladies, you should not speak out
        when you are guessing; but each one of you should secretly write down the
        solutions for me to wrap them up, and take them all in together to await her
        Majesty's personal inspection as to whether they be correct or not."}
}
{
}

\Verse{64}
{
    \zA{}
}
{
    \eA{Upon listening to these words, Pao-ch'ai drew near, and perceived at a
        glance, that it consisted of a stanza of four lines, with seven characters
        in each; but though there was no novelty or remarkable feature about it, she
        felt constrained to outwardly give utterance to words of praise. "It's hard
        to guess!" she simply added, while she pretended to be plunged in thought,
        for the fact is that as soon as she had cast her eye upon it, she had at
        once solved it. Pao-yü, Tai-yü, Hsiang-yün, and T'an-ch'un, had all four
        also hit upon the answer, and each had secretly put it in writing;}
}
{
}

\Verse{65}
{
    \zA{}
}
{
    \eA{and Chia Huan, Chia Lan and the others were at the same time sent for, and
        every one of them set to work to exert the energies of his mind, and, when
        they arrived at a guess, they noted it down on paper; after which every
        individual member of the family made a choice of some object, and composed a
        riddle, which was transcribed in a large round hand, and affixed on the
        lantern. This done, the eunuch took his departure, and when evening drew
        near, he came out and delivered the commands of the imperial consort. "The
        conundrum," he said, "written by Her Highness, the other day, has been
        solved by every one, with the exception of Miss Secunda and master Tertius,
        who made a wrong guess. Those composed by you, young ladies, have likewise
        all been guessed;}
}
{
}

\Verse{66}
{
    \zA{}
}
{
    \eA{but Her Majesty does not know whether her solutions are right or not." While
        speaking, he again produced the riddles, which had been written by them,
        among which were those which had been solved, as well as those which had not
        been solved; and the eunuch, in like manner, took the presents, conferred by
        the imperial consort, and handed them over to those who had guessed right.
        To each person was assigned a bamboo vase, inscribed with verses, which had
        been manufactured for palace use, as well as articles of bamboo for tea;
        with the exception of Ying-ch'un and Chia Huan, who were the only two
        persons who did not receive any. But as Ying-ch'un looked upon the whole
        thing as a joke and a trifle, she did not trouble her mind on that score,
        but Chia Huan at once felt very disconsolate.}
}
{
}

\Verse{67}
{
    \zA{}
}
{
    \eA{"This one devised by Mr. Tertius," the eunuch was further heard to say, "is
        not properly done; and as Her Majesty herself has been unable to guess it
        she commanded me to bring it back, and ask Mr. Tertius what it is about."
        After the party had listened to these words, they all pressed forward to see
        what had been written. The burden of it was this:  The elder brother has
        horns only eight; The second brother has horns only two; The elder brother
        on the bed doth sit; Inside the room the second likes to squat. After
        perusal of these lines, they broke out, with one voice, into a loud fit of
        laughter; and Chia Huan had to explain to the eunuch that the one was a
        pillow, and the other the head of an animal. Having committed the
        explanation to memory and accepted a cup of tea, the eunuch took his
        departure;}
}
{
}

\Verse{68}
{
    \zA{}
}
{
    \eA{and old lady Chia, noticing in what buoyant spirits Yüan Ch'un was, felt
        herself so much the more elated, that issuing forthwith directions to
        devise, with every despatch, a small but ingenious lantern of fine texture
        in the shape of a screen, and put it in the Hall, she bade each of her
        grandchildren secretly compose a conundrum, copy it out clean, and affix it
        on the frame of the lantern; and she had subsequently scented tea and fine
        fruits, as well as every kind of nicknacks, got ready, as prizes for those
        who guessed right. And when Chia Cheng came from court and found the old
        lady in such high glee he also came over in the evening, as the season was
        furthermore holiday time, to avail himself of her good cheer to reap some
        enjoyment.}
}
{
}

\Verse{69}
{
    \zA{}
}
{
    \eA{In the upper part of the room seated themselves, at one table dowager lady
        Chia, Chia Cheng, and Pao-yü; madame Wang, Pao-ch'ai, Tai-yü, Hsiang-yün sat
        round another table, and Ying-ch'un, Tan-ch'un and Hsi Ch'un the three of
        them, occupied a separate table, and both these tables were laid in the
        lower part, while below, all over the floor, stood matrons and waiting-maids
        for Li Kung-ts'ai and Hsi-feng were both seated in the inner section of the
        Hall, at another table.}
}
{
}

\Verse{70}
{
    \zA{}
}
{
    \eA{Chia Chen failed to see Chia Lan, and he therefore inquired: "How is it I
        don't see brother Lan," whereupon the female servants, standing below,
        hastily entered the inner room and made inquiries of widow Li. "He says,"
        Mrs. Li stood up and rejoined with a smile, "that as your master didn't go
        just then to ask him round, he has no wish to come!" and when a matron
        delivered the reply to Chia Cheng; the whole company exclaimed much amused:
        "How obstinate and perverse his natural disposition is!" But Chia Cheng lost
        no time in sending Chia Huan, together with two matrons, to fetch Chia Lan;
        and, on his arrival, dowager lady Chia bade him sit by her side, and, taking
        a handful of fruits, she gave them to him to eat; after which the party
        chatted, laughed, and enjoyed themselves.}
}
{
}

\Verse{71}
{
    \zA{}
}
{
    \eA{Ordinarily, there was no one but Pao-yü to say much or talk at any length,
        but on this day, with Chia Cheng present, his remarks were limited to
        assents. And as to the rest, Hsiang-yün had, though a young girl, and of
        delicate physique, nevertheless ever been very fond of talking and
        discussing; but, on this instance, Chia Cheng was at the feast, so that she
        also held her tongue and restrained her words. As for Tai-yü she was
        naturally peevish and listless, and not very much inclined to indulge in
        conversation; while Pao-ch'ai, who had never been reckless in her words or
        frivolous in her deportment, likewise behaved on the present occasion in her
        usual dignified manner.}
}
{
}

\Verse{72}
{
    \zA{}
}
{
    \eA{Hence it was that this banquet, although a family party, given for the sake
        of relaxation, assumed contrariwise an appearance of restraint, and as old
        lady Chia was herself too well aware that it was to be ascribed to the
        presence of Chia Cheng alone, she therefore, after the wine had gone round
        three times, forthwith hurried off Chia Cheng to retire to rest. No less
        cognisant was Chia Cheng himself that the old lady's motives in packing him
        off were to afford a favourable opportunity to the young ladies and young
        men to enjoy themselves, and that is why, forcing a smile, he observed:
        "Having to-day heard that your venerable ladyship had got up in here a large
        assortment of excellent riddles, on the occasion of the spring festival of
        lanterns, I too consequently prepared prizes, as well as a banquet, and came
        with the express purpose of joining the company;}
}
{
}

\Verse{73}
{
    \zA{}
}
{
    \eA{and why don't you in some way confer a fraction of the fond love, which you
        cherish for your grandsons and granddaughters, upon me also, your son?"
        "When you're here," old lady Chia replied smilingly, "they won't venture to
        chat or laugh; and unless you go, you'll really fill me with intense
        dejection! But if you feel inclined to guess conundrums, well, I'll tell you
        one for you to solve; but if you don't guess right, mind, you'll be
        mulcted!" "Of course I'll submit to the penalty," Chia Cheng rejoined
        eagerly, as he laughed, "but if I do guess right, I must in like manner
        receive a reward!" "This goes without saying!" dowager lady Chia added;
        whereupon she went on to recite:  The monkey's body gently rests on the tree
        top! "This refers," she said, "to the name of a fruit."}
}
{
}

\Verse{74}
{
    \zA{}
}
{
    \eA{Chia Cheng was already aware that it was a lichee, but he designedly made a
        few guesses at random, and was fined several things; but he subsequently
        gave, at length, the right answer, and also obtained a present from her
        ladyship. In due course he too set forth this conundrum for old lady Chia to
        guess:  Correct its body is in appearance,  Both firm and solid is it in
        substance; To words, it is true, it cannot give vent,  But spoken to, it
        always does assent. When he had done reciting it, he communicated the answer
        in an undertone to Pao-yü; and Pao-yü fathoming what his intention was,
        gently too told his grandmother Chia, and her ladyship finding, after some
        reflection, that there was really no mistake about it, readily remarked that
        it was an inkslab. "After all," Chia Cheng smiled;}
}
{
}

\Verse{75}
{
    \zA{}
}
{
    \eA{"Your venerable ladyship it is who can hit the right answer with one guess!"
        and turning his head round, "Be quick," he cried, "and bring the prizes and
        present them!" whereupon the married women and waiting-maids below assented
        with one voice, and they simultaneously handed up the large trays and small
        boxes. Old lady Chia passed the things, one by one, under inspection; and
        finding that they consisted of various kinds of articles, novel and
        ingenious, of use and of ornament, in vogue during the lantern festival, her
        heart was so deeply elated that with alacrity she shouted, "Pour a glass of
        wine for your master!" Pao-yü took hold of the decanter, while Ying Ch'un
        presented the cup of wine. "Look on that screen!"}
}
{
}

\Verse{76}
{
    \zA{}
}
{
    \eA{continued dowager lady Chia, "all those riddles have been written by the
        young ladies; so go and guess them for my benefit!" Chia Cheng signified his
        obedience, and rising and walking up to the front of the screen, he noticed
        the first riddle, which was one composed by the Imperial consort Yüan, in
        this strain:  The pluck of devils to repress in influence it abounds,  Like
        bound silk is its frame, and like thunder its breath resounds. But one
        report rattles, and men are lo! in fear and dread;}
}
{
}

\Verse{77}
{
    \zA{}
}
{
    \eA{Transformed to ashes 'tis what time to see you turn the head. "Is this a
        cracker?" Chia Cheng inquired. "It is," Pao-yü assented. Chia Cheng then
        went on to peruse that of Ying-Ch'un's, which referred to an article of use:
        Exhaustless is the principle of heavenly calculations and of human  skill;
        Skill may exist, but without proper practice the result to find hard  yet
        will be! Whence cometh all this mixed confusion on a day so still? Simply it
        is because the figures Yin and Yang do not agree. "It's an abacus," Chia
        Cheng observed. "Quite so!" replied Ying Ch'un smiling; after which they
        also conned the one below, by T'an-ch'un, which ran thus and had something
        to do with an object:  This is the time when 'neath the stairs the pages
        their heads raise!}
}
{
}

\Verse{78}
{
    \zA{}
}
{
    \eA{The term of "pure brightness" is the meetest time this thing to make! The
        vagrant silk it snaps, and slack, without tension it strays! The East wind
        don't begrudge because its farewell it did take! "It would seem," Chia Cheng
        suggested, "as if that must be a kite!" "It is," answered T'an C'h'un;
        whereupon Chia Cheng read the one below, which was written by Tai-yü to this
        effect and bore upon some thing:  After the audience, his two sleeves who
        brings with fumes replete? Both by the lute and in the quilt, it lacks luck
        to abide! The dawn it marks; reports from cock and man renders effete! At
        midnight, maids no trouble have a new one to provide! The head, it glows
        during the day, as well as in the night!}
}
{
}

\Verse{79}
{
    \zA{}
}
{
    \eA{Its heart, it burns from day to day and 'gain from year to year! Time
        swiftly flies and mete it is that we should hold it dear! Changes might
        come, but it defies wind, rain, days dark or bright! "Isn't this a scented
        stick to show the watch?" Chia Cheng inquired. "Yes!" assented Pao-yü,
        speaking on Tai-yü's behalf; and Chia Cheng thereupon prosecuted the perusal
        of a conundrum, which ran as follows, and referred to an object; With the
        South, it sits face to face,  And the North, the while, it doth face; If the
        figure be sad, it also is sad,  If the figure be glad, it likewise is glad!
        "Splendid! splendid!" exclaimed Chia Cheng, "my guess is that it's a
        looking-glass. It's excellently done!" Pao-yü smiled.}
}
{
}

\Verse{80}
{
    \zA{}
}
{
    \eA{"It is a looking glass!" he rejoined. "This is, however, anonymous; whose
        work is it?" Chia Cheng went on to ask, and dowager lady Chia interposed:
        "This, I fancy, must have been composed by Pao-yü," and Chia Cheng then said
        not a word, but continued reading the following conundrum, which was that
        devised by Pao-ch'ai, on some article or other:  Eyes though it has;
        eyeballs it has none, and empty 'tis inside! The lotus flowers out of the
        water peep, and they with gladness meet,  But when dryandra leaves begin to
        drop, they then part and divide,  For a fond pair they are, but, united,
        winter they cannot greet. When Chia Cheng finished scanning it, he gave way
        to reflection. "This object," he pondered, "must surely be limited in use!}
}
{
}

\Verse{81}
{
    \zA{}
}
{
    \eA{But for persons of tender years to indulge in all this kind of language,
        would seem to be still less propitious; for they cannot, in my views, be any
        of them the sort of people to enjoy happiness and longevity!" When his
        reflections reached this point, he felt the more dejected, and plainly
        betrayed a sad appearance, and all he did was to droop his head and to
        plunge in a brown study. But upon perceiving the frame of mind in which Chia
        Cheng was, dowager lady Chia arrived at the conclusion that he must be
        fatigued; and fearing, on the other hand, that if she detained him, the
        whole party of young ladies would lack the spirit to enjoy themselves, she
        there and then faced Chia Cheng and suggested: "There's no need really for
        you to remain here any longer, and you had better retire to rest;}
}
{
}

\Verse{82}
{
    \zA{}
}
{
    \eA{and let us sit a while longer; after which, we too will break up!" As soon
        as Chia Cheng caught this hint, he speedily assented several consecutive
        yes's; and when he had further done his best to induce old lady Chia to have
        a cup of wine, he eventually withdrew out of the Hall. On his return to his
        bedroom, he could do nothing else than give way to cogitation, and, as he
        turned this and turned that over in his mind, he got still more sad and
        pained. "Amuse yourselves now!" readily exclaimed dowager lady Chia, during
        this while, after seeing Chia Cheng off;}
}
{
}

\Verse{83}
{
    \zA{}
}
{
    \eA{but this remark was barely finished, when she caught sight of Pao-yü run up
        to the lantern screen, and give vent, as he gesticulated with his hands and
        kicked his feet about, to any criticisms that first came to his lips. "In
        this," he remarked, "this line isn't happy; and that one, hasn't been
        suitably solved!" while he behaved just like a monkey, whose fetters had
        been let loose. "Were the whole party after all," hastily ventured Tai-yü,
        "to sit down, as we did a short while back and chat and laugh; wouldn't that
        be more in accordance with good manners?" Lady Feng thereupon egressed from
        the room in the inner end and interposed her remarks. "Such a being as you
        are," she said, "shouldn't surely be allowed by Mr. Chia Cheng, an inch or a
        step from his side, and then you'll be all right.}
}
{
}

\Verse{84}
{
    \zA{}
}
{
    \eA{But just then it slipped my memory, for why didn't I, when your father was
        present, instigate him to bid you compose a rhythmical enigma; and you
        would, I have no doubt, have been up to this moment in a state of
        perspiration!" At these words, Pao-yü lost all patience, and laying hold of
        lady Feng, he hustled her about for a few moments. But old lady Chia went on
        for some time to bandy words with Li Kung-ts'ai, with the whole company of
        young ladies and the rest, so that she, in fact, felt considerably tired and
        worn out;}
}
{
}

\Verse{85}
{
    \zA{}
}
{
    \eA{and when she heard that the fourth watch had already drawn nigh, she
        consequently issued directions that the eatables should be cleared away and
        given to the crowd of servants, and suggested, as she readily rose to her
        feet, "Let us go and rest! for the next day is also a feast, and we must get
        up at an early hour; and to-morrow evening we can enjoy ourselves again!"
        whereupon the whole company dispersed. But now, reader, listen to the sequel
        given in the chapter which follows.}
}
{
}
